% !TEX TS-program = pdflatex
% !TEX encoding = UTF-8 Unicode

% This is a simple template for a LaTeX document using the "article" class.
% See "book", "report", "letter" for other types of document.

\documentclass[11pt]{article} % use larger type; default would be 10pt

\usepackage[utf8]{inputenc} % set input encoding (not needed with XeLaTeX)

%%% Examples of Article customizations
% These packages are optional, depending whether you want the features they provide.
% See the LaTeX Companion or other references for full information.

%%% PAGE DIMENSIONS
\usepackage{geometry} % to change the page dimensions
\geometry{a4paper} % or letterpaper (US) or a5paper or....
% \geometry{margin=2in} % for example, change the margins to 2 inches all round
% \geometry{landscape} % set up the page for landscape
%   read geometry.pdf for detailed page layout information

\usepackage{graphicx} % support the \includegraphics command and options

% \usepackage[parfill]{parskip} % Activate to begin paragraphs with an empty line rather than an indent

%%% PACKAGES
\usepackage{booktabs} % for much better looking tables
\usepackage{array} % for better arrays (eg matrices) in maths
\usepackage{paralist} % very flexible & customisable lists (eg. enumerate/itemize, etc.)
\usepackage{verbatim} % adds environment for commenting out blocks of text & for better verbatim
\usepackage{subfig} % make it possible to include more than one captioned figure/table in a single float
% These packages are all incorporated in the memoir class to one degree or another...

%%% HEADERS & FOOTERS
\usepackage{fancyhdr} % This should be set AFTER setting up the page geometry
\pagestyle{fancy} % options: empty , plain , fancy
\renewcommand{\headrulewidth}{0pt} % customise the layout...
\lhead{}\chead{}\rhead{}
\lfoot{}\cfoot{\thepage}\rfoot{}

%%% SECTION TITLE APPEARANCE
\usepackage{sectsty}
\allsectionsfont{\sffamily\mdseries\upshape} % (See the fntguide.pdf for font help)
% (This matches ConTeXt defaults)

%%% ToC (table of contents) APPEARANCE
\usepackage[nottoc,notlof,notlot]{tocbibind} % Put the bibliography in the ToC
\usepackage[titles,subfigure]{tocloft} % Alter the style of the Table of Contents
\renewcommand{\cftsecfont}{\rmfamily\mdseries\upshape}
\renewcommand{\cftsecpagefont}{\rmfamily\mdseries\upshape} % No bold!

%%% END Article customizations

%%% The "real" document content comes below...

\title{3.1 - Diferencias finitas para derivadas de primer y segundo orden}
\author{Karla Munguia A01741255 \and  Ivor Question}
%\date{} % Activate to display a given date or no date (if empty),
         % otherwise the current date is printed 

\begin{document}
\maketitle

\section{First section}

Sea $f(x)$ una función cuya derivada se desea aproximar por diferencias finitas. Considere la expansión en serie de Taylor de 
$g(\Delta x) = f(x+\Delta x) - f(x-\Delta x)$ como función de $\Delta x$ alrededor de $\Delta x = 0$. 
\begin{enumerate}
\item Muestre que el error de truncamiento para la aproximación de la primera derivada por diferencias centradas es 
$-(\Delta x)^2 f'''(\epsilon_3)/6$.
\item Muestre explícitamente que dicha diferencia centrada es exacta para cualquier polinomio cuadrático.
\end{enumerate}


Derive la fórmula para la aproximación de la segunda derivada de $f(x)$ por diferencias centradas usando dos veces la aproximación por diferencias centradas para la primera derivada.


Supongamos que no conocemos la fórmula de diferencias centradas para aproximar $f''(x)$, pero que es posible encontrar dicha aproximación mediante una combinación lineal de los valores de $f(x_0), f(x_0 - \Delta x)$, y $f(x_0 + \Delta x)$:
$$
\frac{d^2f}{dx^2} \approx af(x_0 - \Delta x) + bf(x_0) + cf(x_0 + \Delta x)
$$
Determine $a,b$ y $c$ expandiendo el lado derecho en serie de Taylor alrededor de $x_0$ e igualando los coeficientes.


\end{document}
